\documentclass[10pt,a4paper]{amsart}
\usepackage[margin=19mm]{geometry}
\usepackage{amsmath,amssymb,mathtools}
\usepackage[scr=rsfs,cal=euler]{mathalfa}
\usepackage{xfrac}
\usepackage[unicode,hidelinks,bookmarks=false]{hyperref}
\newcommand{\ie}{i.e.\ }
\newcommand{\cf}{cf.\ }
\newcommand{\resp}{respectively\ }
\newcommand{\Subsection}{\subsection}
\providecommand{\nobreakdash}{\nobreak\hbox{-}}
\setlength{\emergencystretch}{2em}
\multlinegap0pt
\usepackage{amssymb}
\newtheorem{lem}{Lemma}
\newcommand{\E}{{\bf{E}}}
\newcommand{\PP}{{\bf{P}}}
\title{On the strength of connectedness of unions of random graphs}
\author{Mindaugas Bloznelis}
\date{Source: arXiv:2602.02166v2 (CC BY 4.0)}
\begin{document}
\maketitle

\noindent\emph{Source and licence.} On the strength of connectedness of unions of random graphs. Version arXiv:2602.02166v2 (CC BY 4.0), distributed by arXiv under the Creative Commons Attribution 4.0 International licence, http://creativecommons.org/licenses/by/4.0/, as recorded in the arXiv licence metadata for this exact version and shown on the abstract page https://arxiv.org/abs/2602.02166v2. Full licence notice: \texttt{licenses/arxiv-2602.02166-CC-BY-4.0.txt}.\par\smallskip

\noindent\textit{Editorial scope.} Counted unit: the article's lem AB+ (source0.tex lines 2768--2799). The article's complete original proof of it is delivered with it (source0.tex lines 2801--3561, one \texttt{proof} environment of 761 lines). No mathematics is written, repaired or re-derived: the statement and the proof are the article's own lines, in the article's own notation, and no retained line is substituted or re-flowed.  Retained uncounted as the source-local context the proof needs: lines 1622--1625; lines 2615--2650; lines 2652--2657; lines 2672--2684; lines 2689--2698.  Nothing was omitted from any retained span, so the package carries no omission record.  No retained line contains a citation, so the wrapper carries no bibliography and no bibliography entry.  No retained line contains a figure environment, an image-inclusion command or drawing code, so no image is delivered and the package declares no asset dependency.  Editorial formatting only, in addition: an AMS \texttt{amsart} preamble that restates the article's own declarations and macros used by the retained lines, namely the environments \texttt{lem} on separate counters, together with the standard packages those lines need.  The retained environments print in this document as Lemma 1 in order of appearance, whereas the article prints the same items with the numbering of its own full document.  The complete native source of the article as distributed in the arXiv e-print for this exact version is bundled unchanged as \texttt{sources/193831/source0.tex}.
\begin{align}
\label{condition_X^2}
\limsup_n\E X_{n,i_*}^2<\infty.
\end{align}

\begin{align}
\label{T}
T
&
=
\prod_{i\in[m]}
\PP\{d_i(1)=0\}
=
\prod_{i\in[m]}
\left(1-\frac{x_i}{n}\right)
=
e^{-\frac{m}{n}\kappa}(1+o(1)),
\\
\nonumber
H
&
=
\prod_{i\in[m]}
\PP\{d_i(1)=d_i(2)=0\}
=
\prod_{i\in[m]}
\E
\left(
\left(1-\frac{X_i}{n}\right)
\left(1-\frac{X_i}{n-1}\right)
\right)
\\
\label{H}
&
=
\prod_{i\in[m]}
\left(
1-2\frac{x_i}{n}+\frac{z_i}{(n)_2}\right)
=
e^{-2\frac{m}{n}\kappa}(1+o(1)).
\end{align}

\begin{align}
\label{2026-01-10+6}
\max_{i\in [m]}x_i^2=O(n\ln n),
\qquad 
\max_{i\in [m]}z_i=O(n\ln n).
\end{align}

\begin{align}
\label{2026-01-10+1}
&
\max_{B\subset [m],\, |B|\le b} \prod_{i\in B}\left(1-\frac{x_i}{n}\right)
=
1-O\left(\frac{\sqrt{\ln n}}{\sqrt{n}}\right),
\\
\label{2026-01-10+2}
&
\max_{B\subset [m],\, |B|\le b} \prod_{i\in B}\left(1-2\frac{x_i}{n}+\frac{z_i}{(n)_2}\right)
=
1-O\left(\frac{\sqrt{\ln n}}{\sqrt{n}}\right).
\end{align}

\begin{align}
\label{Fact1}
b!\sum_{B\in\binom{[m]}{b}}\prod_{i\in B}a_i
=
(a_1+\cdots+a_m)^b
-
R
\le 
(a_1+\cdots+a_m)^b,
\end{align}

\begin{lem}\label{AB+} 
Let $a\ge 1$ and $k\ge 0$ be integers.
Let $m,n\to+\infty$.
 Assume that $m=m(n)=o(n\ln^2n)$ and $n\ln n=O(m)$.  Assume that (\ref{condition_X^2}) holds. Then
 \begin{align}
  \label{2025-09-04}
 \PP\{d'(1)=k\}
 &
 =
 \frac{\kappa^k}{k!}\left(\frac{m}{n}\right)^k e^{-\kappa\frac{m}{n}}
 (1+o(1)),
 \\
 \label{2025-06-21_6+r+r}
 \PP\{d_*'(1)=d'(1)=k\}
 &
 =
 \frac{\kappa_a^k}{k!}
 \left(\frac{m}{n}\right)^k
  e^{-\kappa\frac{m}{n}}
 (1+o(1)).
 \end{align}
 Assume, in addition, that $\liminf\kappa_a>0$.
 For $k\ge 0$ we have 
 \begin{align}
  \label{2025-6-21_9+r}
 \PP\{d_*'(1)=d'(1)=k,\, d_*'(2)=d'(2)=k\}
& =\frac{\kappa_a^{2k}}{(k!)^2}
\left(\frac{m}{n}\right)^{2k}
 e^{-2\kappa\frac{m}{n}}
 (1+o(1)).
 \end{align}
\end{lem}

\begin{proof}[ Proof of Lemma \ref{AB+}]
For $k=0$ (\ref{2025-09-04}), (\ref{2025-06-21_6+r+r}) follow from (\ref{T}) via identities
\begin{align*}
\PP\{d'(1)=0\}
=
\PP\{d'(1)=d'_*(1)=0\}
= 
T.
\end{align*}
Similarly,  (\ref{2025-6-21_9+r}) follows from (\ref{H})
via identity
\begin{align*}
\PP\{d_*'(1)=d'(1)=0,\, d_*'(2)=d'(2)=0\}
=
\PP\{d'(1)=d'(2)=0\}=H.
\end{align*}

For the rest of the proof we assume that $k\ge 1$.



Proof of (\ref{2025-09-04}).
 We apply the total probability formula and use the independence of 
 $G_1,\dots, G_m$:
 \begin{align*}
 \PP\{d'(1)=k\}
&
=
\sum_{B\in\binom{[m]}{k}}
\PP
\left\{
d_i(1)>0,\, i\in B
\
\ 
{\text{and}}
\
\
d_j(1)=0, \, j \in[m]\setminus B
\right\}
\\
\nonumber
&
=
\sum_{B\in\binom{[m]}{k}}
\left(
\prod_{i\in B}
\PP\{d_i(1)>0\}
\right)
\left(
\prod_{i\in[m]\setminus B}
\PP\{d_j(1)=0\}
\right)
\\
\nonumber
&=
\sum_{B\in\binom{[m]}{k}}
\left(
\prod_{i\in B}
\frac{x_{i}}{n}
\right)
\left(
\prod_{j\in[m]\setminus B}
\left(
1
-
\frac{x_j}{n}
\right)
\right)
\\
&
=
T
\sum_{B\in\binom{[m]}{k}}
\prod_{i\in B}\frac{x_{i}}{n}
\left( 1-\frac{x_i}{n}\right)^{-1}.
\end{align*}
In view of (\ref{2026-01-10+1}) we 
have $\prod_{i\in B}
\left( 1-\frac{x_i}{n}\right)^{-1}=
1+o(1)$ uniformly 
in  $B$.
Hence,
 \begin{align*}
 \PP\{d'(1)=k\}
=
(1+o(1))
T\frac{1}{n^k}
\sum_{B\in\binom{[m]}{k}}
\prod_{i\in B}x_{i}
\end{align*}
Finally, 
invoking 
 (\ref{T}) and approximation 
$\sum_{B\in\binom{[m]}{k}}
\prod_{i\in B}x_{i}
=
\frac{m^k}{k!}\kappa^k+O(m^{k-1})$, which follows by 
(\ref{Fact1}),  we obtain 
 (\ref{2025-09-04}).

Proof of (\ref{2025-06-21_6+r+r}). We proceed similarly as in the proof of (\ref{2025-09-04})
above. We have
\begin{align*}
\PP\{d'_*(1)=d'(1)=k\}
&
=
\sum_{B\in\binom{[m]}{k}}
\PP\{d_i(1)=a,\, i\in B
\
\
{\text{and}}
\
\
d_j(1)=0,
\, j\in [m]\setminus B\}
\\
\nonumber
&
=
\sum_{B\in\binom{[m]}{k}}
\left(
\prod_{i\in B}
\PP\{d_i(1)=a\}
\right)
\left(
\prod_{i\in[m]\setminus B}
\PP\{d_j(1)=0\}
\right)
\\
\nonumber
&=
\sum_{B\in\binom{[m]}{k}}
\left(
\prod_{i\in B}
\frac{x_{a,i}}{n}
\right)
\left(
\prod_{j\in[m]\setminus B}
\left(
1
-
\frac{x_j}{n}
\right)
\right)
%\\
%&
%=
%T
%\sum_{B\in\binom{[m]}{k}}
%\prod_{i\in B}\frac{x_{a,i}}{n}
%\left( 1-\frac{x_i}{n}\right)^{-1}
\\
\nonumber
&
=
(1+o(1))
T\frac{1}{n^k}
\sum_{B\in\binom{[m]}{k}}
\prod_{i\in B}x_{a,i}
\\
\nonumber
&
=
 (1+o(1))\frac{m^k}{n^k}\frac{\kappa_a^k}{k!}
 e^{-\frac{m}{n}\kappa}.
\end{align*}

Proof of (\ref{2025-6-21_9+r}). Denote for short
$p=\PP\{d'_*(1)=d'(1)=d'_*(2)=d'(2)=k\}$.
Given mutually disjoint sets $B_1,B_2,B_3\subset [m]$, let 
$I_{B_1,B_2,B_3}$
denote the event that
\begin{align*}
&
d_i(1)=d_i(2)=a
\
 \forall i\in B_1,
\quad
d_j(1)=a, 
\
d_j(2)=0
\
 \forall j\in B_2,
\\
&
d_h(1)=0, 
\ 
d_h(2)=a
\
 \forall h\in B_3,
\quad
d_\ell(1)=d_\ell(2)=0
\
\forall \ell\in [m]\setminus (B_1\cup B_2\cup B_3).
 \end{align*}
By the total probability formula we have 
\begin{align*}
p
=
\sum_{s=0}^k{\bar p}_s,
\qquad
{\bar p}_s
=
\sum_
{\substack{B_1\subset [m]\\ |B_1|=s}}
\
\sum_{\substack{B_2\subset [m]\setminus B_1\\
|B_2|=k-s}}
\
\sum_{\substack{B_3\subset [m]\setminus (B_1\cup B_2)\\
|B_3|=k-s}}
\PP\{I_{B_1,B_2,B_3}\}.
\end{align*} 
Furthermore, by the independence of $G_1,\dots, G_m$,
we factorize the probability 
\begin{align}
\nonumber
\PP\{I_{B_1,B_2,B_3}\}
&=
\left(
\prod_{i\in B_1}
p_i
\right)
\times
\left(
\prod_{j\in B_2}p_j(1,2)
\right)
\times
\left(
\prod_{h\in B_3}p_h(2,1)
\right)
\times
\prod_{\ell\in [m]\setminus(B_1\cup B_2\cup B_3)}
q_\ell
\\
\label{2026-01-10+4}
&
=:
P_1(B_1)
\times
P_2(B_2)
\times
P_3(B_3)
\times
Q(B_1,B_2,B_3).
\end{align}
Here we denote
 \begin{align*}
&
p_i=\PP\{d_i(1)=d_i(2)=a\},
\qquad
q_i=\PP\{d_i(1)=d_i(2)=0\},
\\
&
p_i(1,2)=\PP\{d_i(1)=a, \,d_i(2)=0\},
\qquad
p_i(2,1)=\PP\{d_i(1)=0,\, d_i(2)=a\}.
\end{align*}
A calculation shows that
\begin{align*}
p_i
&
=
\E \frac{(X_{a,i})_2}{(n)_2}=\frac{z_{a,i}}{(n)_2},
\qquad
q_i
=
\E\frac{(n-X_i)_2}{(n)_2}
=1-2\frac{x_i}{n}+\frac{z_i}{(n)_2},
\\
p_i(1,2)
&
=
p_i(2,1)
=
\E\frac{X_{a,i}(n-X_i)}{(n)_2}
=
\frac{x_{a,i}}{n-1}-\frac{\E (X_iX_{a,i})}{(n)_2}
\le 
\frac{x_{a,i}}{n-1}.
\end{align*}
%\end{proof}
%\end{document}
We also note that in view of (\ref{2026-01-10+2})  the last term on the right of (\ref{2026-01-10+4})
\begin{align}
\label{2026-01-12+5}
%\prod_{\ell\in [m]\setminus(B_1\cup B_2\cup B_3)}
%q_\ell
Q(B_1,B_2,B_3)
=
H\prod_{\ell\in B_1\cup B_2\cup B_3}
\left(
1-2\frac{x_i}{n}+\frac{z_i}{(n)_2}
\right)^{-1}
=
H
(1+o(1)),
\end{align}
where the bound $o(1)$ holds
uniformly over $B_1,B_2,B_3$ satisfying $|B_1\cup B_2\cup B_3|\le 2k$.



In the remaining part of the proof we show that
\begin{align}
\label{2026-01-10+5}
{\bar p}_0
=
\frac{\kappa_a^{2k}}{(k!)^2}
\left(
\frac{m}{n}
\right)^{2k}
H
\left(
1
+
o(1)
\right)
\qquad
{\text{and}}
\qquad
\sum_{s=1}^k {\bar p}_s
=
O\left(H\frac{m^{2k-1}}{n^{2k}}\right).
\end{align}
These relations combined with (\ref{H}) imply
(\ref{2025-6-21_9+r}).



Before the proof  of
(\ref{2026-01-10+5}) 
we introduce some notation.
For an integer $b\ge 1$ we denote
\begin{align*}
S_{1,b}
=\sum_
{\substack{B_1\subset [m],\\ |B_1|=b}}
\prod_{i\in B_1}
p_i,
\quad
\
S_{2,b}=
\sum_{\substack{B_2\subset [m]\\
|B_2|=b}}
\prod_{j\in B_2}p_j(1,2),
\quad
\
S_{3,b}=
\sum_{\substack{B_3\subset [m]\\
|B_3|=b}}
\prod_{h\in B_3}p_h(2,1).
\end{align*}
Note that $S_{2,b}=S_{3,b}$.
Next we establish several useful facts about the sums $S_{1,b}$ and $S_{2,b}$.

Using
$p_i=\frac{z_{a,i}}{(n)_2}$ 
and
$p_i(1,2)
=
p_i(2,1)
\le 
\frac{x_{a,i}}{n-1}$, 
and (\ref{Fact1}) we upperbound
\begin{align}
\label{2026-01-12+2}
S_{1,b}
\le 
\frac{1}{b!}
\left(\frac{m}{(n)_2}\right)^b
\left(
\E(X_{a,i_*})_2
\right)^b,
\qquad
S_{r,b}
\le 
\frac{1}{b!}
\left(\frac{m}{n-1}\right)^b
\left(
\E X_{a,i_*}
\right)^b,
\quad
r=2,3.
\end{align}
 Moreover, 
%{\color{red}for $b\ge 2$} 
the second relation can be upgraded
to the approximate identity (recall that $\E X_{a,i_*}=\kappa_a$)
\begin{align}
\label{2026-01-12+3}
S_{2,b}
=
 \frac{\kappa_a^b}{b!}\left(\frac{m}{n}\right)^b
 \left(1+O(n^{-1})\right).
\end{align}
Let us show (\ref{2026-01-12+3}). 
Our conditions $\liminf_n\kappa_a>0$ and (\ref{condition_X^2}) imply 
$\kappa_a^{-1}\E(X_{i_*}X_{a,i_*})=O(1)$.
Now, for $b=1$ we have 
\begin{align*}
S_{2,1}
=
\sum_{j\in [m]}p_j(1,2)
=m\left(\frac{\kappa_a}{n-1}-\frac{\E (X_{i_*}X_{a,i_*})}{(n)_2}\right)
=
\kappa_a
\frac{m}{n}
\left(
1
+
O(n^{-1})
\right).
\end{align*}
For $b\ge 2$ 
 we combine 
 the upper bound 
$S_{2,b}
\le
\frac{\kappa_a^b}{b!}
\left(\frac{m}{n}\right)^b
\left(1+O\left(\frac{1}{n}\right)\right)$,
which 
follows from the the second inequality of 
(\ref{2026-01-12+2}), with a matching lowe bound. We show
the
lower bound using 
(\ref{Fact1}), we have
\begin{align*}
S_{2,b}
\ge \frac{S}{b!}-\frac{R}{2(b-2)!},
\end{align*}
where
\begin{align*}
S=\left(
\sum_{i=1}^mp_i(2,1)\right)^b
\qquad
{\text{and}}
\qquad
R=\left(
\sum_{i=1}^mp_i(2,1)\right)^{b-2}
\sum_{i=1}^mp_i^2(2,1).
\end{align*}
Invoking 
$p_i(1,2)= \frac{x_{a,i}}{n-1}-\frac{\E (X_iX_{a,i})}{(n)_2}$
we write $S$ in the form 
\begin{align*}
S
&
=
\left(\frac{m}{n-1}\right)^b
\left(\kappa_a-\frac{\E(X_{i_*}X_{a,i_*})}{n}\right)^b
=
\left(\frac{m}{n}\right)^b
\kappa_a^b 
\left(1+O(n^{-1})\right).
\end{align*}
Furthermore, using  
$p_i(1,2)\le \frac{x_{a,i}}{n-1}$ 
we upper bound
\begin{align*}
R
\le 
\left(\frac{m}{n-1}\kappa_a\right)^{b-2}
\frac{m}{(n-1)^2}\E X^2_{a,i_*}
=
O\left(\kappa_a^{b-2}\E X_{a,i_*}^2
\frac{m^{b-1}}{n^b}\right)
=
O\left(\kappa_a^b
\frac{m^{b-1}}{n^b}\right).
\end{align*}
In the last step we used the bound
$\kappa_a^{-2}\E X_{a,i_*}^2=O(1)$, which follows from 
our conditions
(\ref{condition_X^2})  and
$\liminf_n\kappa_a>0$.
We arrive to the bound 
$S_{2,b}
\ge 
\frac{\kappa_a^b}{b!}
\left(\frac{m}{n}\right)^b
\left(
1
-
O(n^{-1})
\right)$
thus completing the proof of
(\ref{2026-01-12+3}).

Let us show (\ref{2026-01-10+5}) for $k=1$. 
We have
\begin{align*}
{\bar p}_1
=
\sum_{i\in [m]}p_iQ(\{i\},\emptyset,\emptyset)
=
\sum_{i\in [m]}p_iH(1+o(1))
=
S_{1,1}H(1+o(1))
=
O\left(\frac{m}{n^2}H\right).
\end{align*}
Here we approximated $Q(\{i\},\emptyset,\emptyset)
=
H(1+o(1))$ by  (\ref{2026-01-12+5}) 
and then bounded $S_{1,1}$ by  (\ref{2026-01-12+2}).
We similarly approximate
\begin{align*}
{\bar p}_0
&
=
\sum_{j\in[m]}p_j(1,2)\sum_{h\in[m]\setminus\{j\}}p_h(2,1)Q(\emptyset,\{j\},\{h\})
\\
&
=
\sum_{j\in[m]}p_j(1,2)\sum_{h\in[m]\setminus\{j\}}p_h(2,1)H(1+o(1))
\\
&
=
\left(
S_{2,1}
S_{3,1}
-
\sum_{j\in[m]}p_j(1,2)p_j(2,1)
\right)
H(1+o(1)).
\end{align*}
Invoking the approximation $S_{2,1}S_{3,1}=S_{2,1}^2
=\kappa_a^2(m/n)^2(1+o(1))$, see (\ref{2026-01-12+3}), and bound
\begin{align*}
\sum_{j\in[m]}p_j(1,2)p_j(2,1)
\le \sum_{j\in[m]}\frac{x_{a,j}^2}{(n-1)^2}
=
m\frac{\E X_{a,i_*}^2}{(n-1)^2}
=
O\left(\frac{m}{n^2}\right)
\end{align*}
we obtain 
${\bar p}_0
=
\kappa_a^2(m/n)^2H(1+o(1))+O\left(\frac{m}{n^2}H\right)$
thus arriving to  (\ref{2026-01-10+5}).


Now  we show (\ref{2026-01-10+5}) for $k\ge 2$. 
Let us upper bound ${\bar p}_s$ for $s\ge 1$.
By increasing the range of summation we upper bound
\begin{align*}
{\bar p}_s
&
\le
\sum_
{\substack{B_1\subset [m]\\ |B_1|=s}}
\
\sum_{\substack{B_2\subset [m]\\
|B_2|=k-s}}
\
\sum_{\substack{B_3\subset [m]\\
|B_3|=k-s}}
P_1(B_1)P_2(B_2)P_3(B_3)Q(B_1,B_2,B_3)
\\
&
=
S_{1,s}S_{2,k-s}S_{3,k-s}H(
1+o(1)).
\end{align*}
In the last step we invoked (\ref{2026-01-12+5}).
Now (\ref{2026-01-12+2}) implies
${\bar p}_s=O\left(H\frac{m^{2k-s}}{n^{2k}}\right)$. 
Consequently, we obtain 
$\sum_{s=1}^k{\bar p}_s
=O\left(H\frac{m^{2k-1}}{n^{2k}}\right)$.

Let us show the first relation of (\ref{2026-01-10+5}).
We write ${\bar p}_0$ in the form
${\bar p}_0
=
H{\tilde p}_0+{\tilde R}$,
where
\begin{align*}
{\tilde p}_0
&
=
\sum_{\substack{B_2\subset [m]\\
|B_2|=k}}
P_2(B_2)
\sum_{\substack{B_3\subset [m]\setminus B_2\\
|B_3|=k}}
P_3(B_3),
\\
{\tilde R}
&
=
\sum_{\substack{B_2\subset [m]\\
|B_2|=k}}
P_2(B_2)
\sum_{\substack{B_3\subset [m]\setminus B_2\\
|B_3|=k}}
P_3(B_3) (Q(\emptyset, B_2,B_3)-H).
\end{align*}
Using  (\ref{2026-01-12+5}) we estimate
${\tilde R}=o\left({\tilde p}_0H
\right)$.
Furthermore, using
(\ref{2026-01-12+2}) we estimate
\begin{align*}
{\tilde p}_0
\le 
S_{2,k}S_{3,k}
=
O
\left(\frac{m^{2k}}{n^{2k}}\right).
\end{align*}
We conclude that 
${\tilde R}
=
o\left(H\frac{m^{2k}}{n^{2k}}
\right)$.
Let us evaluate ${\tilde p}_0$.
We have
\begin{align}
\label{2026-01-26}
{\tilde p}_0
&=
\sum_{\substack{B_2\subset [m]\\
|B_2|=k}}
P_2(B_2)
\Biggl(
S_{3,k}
-
\sum_{\substack{B_3\subset [m],|B_3|=k\\
B_3\cap B_2\not=\emptyset}}
P_3(B_3)
\Biggr)
=
S_{2,k}S_{3,k}-{\bar R},
\end{align}
where 
${\bar R}
=
R_1+\dots+R_k$, 
and where
\begin{align*}
R_\ell
=
\sum_{\substack{B_2\subset [m]\\
|B_2|=k}}
P_2(B_2)
{\tilde S}_\ell(B_2),
\qquad
{\tilde S}_\ell(B_2)
=
\sum_{\substack{B_3\subset [m],|B_3|=k\\
|B_3\cap B_2|=\ell}}
P_3(B_3).
\end{align*}
Next, we evaluate the product $S_{2,k}S_{3,k}=S_{2,k}^2$ 
using (\ref{2026-01-12+3}) and show below that
${\bar R}=o
\left(
\left(\frac{m}{n}\right)^{2k-1}
\right)$. Now (\ref{2026-01-26}) implies 
 the first relation of 
(\ref{2026-01-10+5}).






It remains to show that ${\bar R}=o
\left(
\left(\frac{m}{n}\right)^{2k-1}
\right)$. Let us consider the sum
${\tilde S}_\ell(B_2)$.
Given $B_2$, we split $B_3=A\cup D$, 
where $A=B_3\cap B_2$  and $D\cap B_2=\emptyset$.
Clearly, $|A|=\ell$ 
and  $|D|=k-\ell$. 
%We have $P_3(B_3)=P_3(A)P_3(D)$.
Using (\ref{2026-01-10+6}) we upperbound
\begin{align*}
P_3(A)
=
\prod_{h\in A}p_h(2,1)
\le 
\prod_{h\in A}\frac{x_{a,h}}{n-1}
\le 
\prod_{h\in A}\frac{x_{h}}{n-1}
=
O
\left(
\left(
\frac{\sqrt{\ln n}}{\sqrt{n}}
\right)^\ell
\right)
\end{align*}
 uniformly in $A\subset [m]$, $|A|=\ell$. 
 It follows from the identity $P_3(B_3)=P_3(A)P_3(D)$
 that
\begin{align*}
{\tilde S}_\ell(B_2)
&
=
\sum_{\substack{A\subset B_2\\
|A|=\ell}}
P_3(A)
\sum_{\substack{D\subset [m]\setminus B_2\\
|D|=k-\ell}}
P_3(D)
\le
\sum_{\substack{A\subset B_2\\
|A|=\ell}}
P_3(A)
S_{3,k-\ell}
\\
&
\le 
O
\left(
\left(
\frac{\sqrt{\ln n}}{\sqrt{n}}
\right)^\ell
\right)
\binom{k}{\ell}
S_{3,k-\ell}.
\end{align*}
Since the bound hols uniformly over $B_2$, we have
\begin{align*}
R_\ell
=
O
\left(
\left(
\frac{\sqrt{\ln n}}{\sqrt{n}}
\right)^\ell
\right)
\binom{k}{\ell}
S_{3,k-\ell}
S_{2,k}
=
O
\left(
\left(
\frac{\sqrt{\ln n}}{\sqrt{n}}
\right)^\ell
\right)
O\left(\left(\frac{m}{n}\right)^{2k-\ell}
\right).
\end{align*}
In the last step we invoked the upper bounds for $S_{3,k-\ell}$ 
and
$S_{2,k}$ 
shown in (\ref{2026-01-12+2}).
It follows that ${\bar R}=R_1+\cdots+R_k=O
\left(
\frac{\sqrt{\ln n}}{\sqrt{n}}
\left(\frac{m}{n}\right)^{2k-1}
\right)$.
\end{proof}

\end{document}
